\scanpage{46}\mspage{2}

% letterhead
{\footnotesize
\begin{tabular}{@{}l@{\quad}l@{}}
\deleted{MINISTERE DE L'EDUCATION NATIONALE} & \deleted{AVIGNON, le}\\
\deleted{ACADEMIE D'AIX-MARSEILLE} &\\
\hfil\rule{1.5cm}{0.4pt}\hfil &\\
\deleted{Inspection Académique} & \deleted{L'Inspecteur d'académie,}\\
\deleted{de Vaucluse} & \deleted{directeur des services départementaux}\\
\hfil\rule{1.5cm}{0.4pt}\hfil & \deleted{de l'Education Nationale de Vaucluse,}\\
\deleted{Tél. : 90 82 99 76} &\\
\deleted{Référence : Secrétariat n\textsuperscript{o}} & \deleted{à M}
\end{tabular}}
% end letterhead

à l'hypothèse b.

\deleted{OPS} $M$ réalisé comme monoïde d'endomorphismes
d'un ensemble \uline{fini} $E$ (p. ex.
par la représentation régulière gauche de $M$).
La relation $(\times\times)$ implique
\[
\mathrm{Im}\,p\subset\mathrm{Im}\,u\qquad\forall u\in M,
\]
donc \inserted{$E_0=$} $\mathrm{Im}\,p$ est le plus petit des $\mathrm{Im}\,u$,
pour $u\in M$. Soit $u\in M$, on a
\[
u(E_0)=up(E)=\mathrm{Im}\,up)\supset E_0,
\]
\margin{\unsure{?}}et comme $E_0$ est fini, on en conclut que
$u(E_0)=E_0$ et que $u\mid E_0$ est un autom. de
l'ens. $E_0$. Soit $u_{E_0}$ cet autom. L'application
\[
u\longmapsto u_{E_0}\ :\quad M\longrightarrow\mathrm{Aut}(E_0)
\]
est un hom de $M$ dans un groupe. L'hyp.
a) sur $M$ implique donc qu'on a $u_{E_0}=\mathrm{id}_{E_0}$
$\forall u\in M$. On en conclut aussi que
\[
E_0=\text{ens. des pts fixes de }M\text{ opérant dans } E,
\]
car si $x\in E\setminus E_0$, comme $p(x)\in E_0$, \deleted{on} a $p(x)\ne x$,

% archive sheet number: 45
